\documentclass[conference]{IEEEtran}
\usepackage{cite}
\usepackage{amsmath,amssymb,amsfonts}
\usepackage{graphicx}
\usepackage{textcomp}
\usepackage{xcolor}
\usepackage{booktabs}
\usepackage{microtype}
\usepackage{url}
\usepackage{listings}

\lstset{
  basicstyle=\ttfamily\footnotesize,
  breaklines=true,
  frame=single,
  backgroundcolor=\color{gray!10}
}

\begin{document}

\title{Collaborative LaTeX Authoring and Version Control with Overleaf: A Practical Guide for the SmartLuggage IoT Project}

\author{
    \IEEEauthorblockN{Jorge Camacho Camacho, Teammate 2, and Teammate 3}
    \IEEEauthorblockA{\textit{Department of Information and Communication Technologies} \\
    \textit{Universitat Pompeu Fabra (UPF)}, Barcelona, Spain}
}

\maketitle

\begin{abstract}
Writing scientific technical reports for Internet of Things (IoT) systems requires seamless synchronization among multi-author teams, precise mathematical formatting, consistent bibliography referencing, and version-controlled asset pipelines. This document provides a complete guide on utilizing Overleaf in conjunction with GitHub for the \textbf{AeroTrack IoT} project. It details project import workflows, live real-time editing conventions, BibTeX reference management, vector graphics integration, and bi-directional GitHub synchronization.
\end{abstract}

\begin{IEEEkeywords}
Overleaf, LaTeX, Collaborative Writing, Git Integration, Technical Report, IEEEtran.
\end{IEEEkeywords}

\section{Introduction to Overleaf}
Overleaf is a cloud-based LaTeX collaborative editor that offers real-time Google Docs-style co-authoring while retaining the typesetting power of full TeX Live distributions. For teams of three engineering students, Overleaf eliminates the classic \textit{"it doesn't compile on my machine"} problem, as compiler environments and packages are standardized in the cloud.

\section{Project Import and Setup}
There are two fast methods to initialize the AeroTrack IoT paper on Overleaf:

\subsection{Method 1: Direct Zip Upload (Fastest, 30 seconds)}
\begin{enumerate}
    \item Navigate to your Overleaf dashboard at \url{https://www.overleaf.com}.
    \item Click on the green \textbf{New Project} button in the top left.
    \item Select \textbf{Upload Project}.
    \item Drag and drop the \texttt{paper.zip} file from this repository.
    \item Overleaf will unpack all files (\texttt{main.tex}, \texttt{IEEEtran.cls}, \texttt{references.bib}, and the \texttt{figures/} folder) and immediately compile the document.
\end{enumerate}

\subsection{Method 2: Direct GitHub Synchronization}
If you have an Overleaf account linked to GitHub:
\begin{enumerate}
    \item In Overleaf, click \textbf{New Project} $\rightarrow$ \textbf{Import from GitHub}.
    \item Choose the repository: \texttt{joorgecamacho/smart-luggage-iot}.
    \item Overleaf creates a connected clone. Every push to GitHub can be pulled into Overleaf, and edits made in Overleaf can be committed directly to GitHub via Overleaf's \textbf{Git} menu.
\end{enumerate}

\section{Repository Structure Explained}
The paper repository is organized according to academic best practices:
\begin{itemize}
    \item \texttt{IEEEtran.cls}: Official IEEE conference LaTeX document class defining two-column typography, margins, and header styles.
    \item \texttt{main.tex}: Master document containing all mandatory report sections matching the project evaluation rubric.
    \item \texttt{references.bib}: BibTeX database storing structured citations for papers, RFCs, and hardware datasheets.
    \item \texttt{figures/}: High-resolution vector PDF and PNG diagrams. Always reference figures with relative paths (e.g., \texttt{figures/fig1\_system\_architecture.pdf}).
\end{itemize}

\section{Collaborative Workflow Rules}
To ensure efficient pair writing among the three team members:
\begin{enumerate}
    \item \textbf{Section Ownership}: Divide sections by expertise (e.g., Author 1: System Architecture and Protocol; Author 2: App Design and State Machine; Author 3: Demo, Code, and Results).
    \item \textbf{Review Mode}: Use Overleaf's \textit{Track Changes} and \textit{Comments} features on the right-hand margin to discuss revisions without deleting text prematurely.
    \item \textbf{Citations}: Never hardcode numeric citation brackets \texttt{[1]}. Always insert citations using \texttt{\textbackslash cite\{bibkey\}}, referencing entries defined in \texttt{references.bib}.
    \item \textbf{Cross-referencing}: Reference figures, tables, and sections dynamically using \texttt{\textbackslash ref\{fig:label\}} and \texttt{\textbackslash pageref\{...\}}.
\end{enumerate}

\section{Handling LaTeX Errors and Compilation}
Overleaf provides instant syntax highlighting. Common pitfalls and solutions include:
\begin{itemize}
    \item \textit{Special characters}: Symbols like \texttt{\%}, \texttt{\&}, \texttt{\_}, and \texttt{\$} must be escaped (\texttt{\textbackslash\%}, \texttt{\textbackslash\_}) unless inside math mode or listings.
    \item \textit{Missing package errors}: All common packages (\texttt{amsmath}, \texttt{graphicx}, \texttt{booktabs}, \texttt{listings}) are pre-loaded in \texttt{main.tex}.
    \item \textit{Cache clearing}: If compilation behaves unexpectedly after file renames, click the small trash icon next to \textit{Recompile} to clear cached auxiliary files.
\end{itemize}

\section{Exporting the Final Deliverable}
Once finalized:
\begin{enumerate}
    \item Click the \textbf{Download PDF} button next to the preview pane.
    \item Ensure the generated PDF matches the course page length and layout constraints.
    \item Commit your updated \texttt{main.tex} back to the GitHub repository.
\end{enumerate}

\end{document}
